\subsection{自反空间}

定义 3. --- 若典范映射 \(c_{\mathrm{E}}\)（从 E 到 \(\mathbf{E}''\) 中的）是从 E 到 \(\mathbf{E}_b'\) 的强对偶上的一个拓扑向量空间同构，则称局部凸空间 E 是自反的。

特别地，自反空间是半自反的，从而是 Hausdorff 的。

命题 4. --- 自反空间的强对偶是自反的。

这立即由定义 3 得出。

定理 2. --- 为使局部凸 Hausdorff 空间 E 是自反的，必要且充分的是它是桶型的，且 E 的每个有界子集对于弱化拓扑 \(\sigma(E, E')\) 相对紧。

由（IV，第 15 页的）定理 1，这等价于说 E 是自反的当且仅当它是半自反且桶型的。

若 E 是自反的，则 \(E_{b}^{\prime}\) 是自反的（命题 4），从而 E 是桶型的（IV，第 15 页，定理 1）。反之，若 E 是半自反且桶型的，则 \(c_{E}\) 是双射且由 IV 第 15 页的命题 2 是双方连续的，故 E 是自反的。

注. --- * 1) 设 E 是一个无限维实 hilbertian 空间。设 F 表示赋以弱化拓扑的空间 E。空间 E 与 F 具有相同的对偶 E'，且 E 是自反 Banach 空间（V，第 17 页）。因此，F 是半自反的。

然而，E 上的强拓扑与弱化拓扑是不同的，故 F 不是自反的。*

\begin{enumerate}
\def\labelenumi{\arabic{enumi})}
\setcounter{enumi}{1}
\tightlist
\item
  设 E 是一个自反空间，M 是 E 的一个闭向量子空间。M 与 E/M 可能都不是自反空间（IV，第 63 页，习题 10）。* 关于赋范空间的情形，见 IV 第 17 页的命题 7。*
\end{enumerate}

